% Page 081

\seite{81}

Der Abschluß des alten Jahres und der Beginn des\ob
neuen waren für Schule und Gemeinde wenig erfreu\ohb
lich. Denn Lehrer Philipp wurde am 7.\ Dezember 1926\ob
infolge der Komplikationen\unsicher, so daß er ein\ob
Vierteljahr lang an schwerer, hartnäckiger Gaumen\ohb
entzündung und Mittelohrentzündung litt. Der\ob
Unterricht wurde infolgedessen nicht\unsicher{} abgehalten, die\ob
letzten Wochen vor Ostern — zum\unsicher{} — am\ob
13.\ Dezember. Halbtag Dienstag — Nachmittags von 1–5 Uhr\ob
einsetzen.

Ab 24.\ Februar übernahm der Pfälzische Lehrerverbands\ohb
lehrer Emil Frantz aus Trier — von der Regierung geschickt —\ob
den Unterricht in voller Stundenzahl. Damit fand seit Wochen\ob
das geregelte Schulleben\unsicher.

10.\ Mai.    Die Osterferien fielen in die Zeit vom 13.\ bis 25.\ April.\ob
Lehrer Philipp übernahm mit dem heutigen Tag\ob
wieder den Unterricht, frisch und gesundheitlich\unsicher\ob
zum Schuldienst.

Bei im vergangenen Herbst von der Freiwilli\ohb
gen Feuerwehr zur Einladung zum Brandmeister-Kur\ohb
sus hatten neben fast allen jüngern Männern teil\ohb
genommen\unsicher. Die\unsicher{} Feuerwehr\unsicher{} des\ob
Herrn Brand\unsicher{} Herr Lehmann\unsicher{} gewählt.\unsicher\ob
in der\unsicher{} Branddirektors\unsicher{} Peter Simon\unsicher\ob
Velten.\unsicher

4.\ Juni     Mit heutigem Tag beginnen die Pfingstferien. Die\ob
endigen mit dem 13.\ Juni.\ob
Der Stand der Feldfrüchte ist ein normaler. Der\ob
Getreidestand ist gut.\unsicher{} der Nachtfrost\unsicher, ward der\ob
\unsicher Teil mit die Ernte\unsicher.

13.\ Juli    Die Hundstags-Ferien vom 1.– 14.\ Juli.\ob
Gestern Nachmittag gab es vom 5 Uhr so ein schwer\ohb
es Gewitter\unsicher.
\pend
